
\chapter{2021年上海卷（秋）}

\section{填空题}

\begin{problem}
已知 \(z_1=1+\i\)，\(z_2=2+3\i\)，求 \(z_1+z_2=\)\fillinblank{}.
\end{problem}

\begin{answer}
\(3+4\i\)
\end{answer}

\begin{solution}
因为 \(z_1=1+\i\)，\(z_2=2+3\i\)，所以 \(z_1+z_2=3+4\i\). 故答案为 \(3+4\i\).
\end{solution}

\begin{problem}
已知 \(A=\{x\mid 2x\le 1\}\)，\(B=\{-1,0,1\}\)，则 \(A\cap B=\)\fillinblank{}.
\end{problem}

\begin{answer}
\(\{-1,0\}\)
\end{answer}

\begin{solution}
因为 \(A=\{x\mid 2x\le 1\}=\left\{x\mid x\le \frac12\right\}\)，\(B=\{-1,0,1\}\)，所以 \(A\cap B=\{-1,0\}\). 故答案为 \(\{-1,0\}\).
\end{solution}

\begin{problem}
若 \(x^2+y^2-2x-4y=0\)，求圆心坐标为\fillinblank{}.
\end{problem}

\begin{answer}
\((1,2)\)
\end{answer}

\begin{solution}
由 \(x^2+y^2-2x-4y=0\) 可得圆的标准方程为 \(\left( x-1 \right)^2+\left( y-2 \right)^2=5\)，所以圆心坐标为 \((1,2)\). 故答案为 \((1,2)\).
\end{solution}

\begin{problem}
如图正方形 \(ABCD\) 的边长为 \(3\)，求 \(\overrightarrow{AB}\cdot \overrightarrow{AC}=\)\fillinblank{}.

\bitmapfigure[max width=6.0cm]{img/2021/shanghai/q4_fig1.png}
\end{problem}

\begin{answer}
\(9\)
\end{answer}

\begin{solution}
由数量积的定义，可得 \(\overrightarrow{AB}\cdot \overrightarrow{AC}=\left| \overrightarrow{AB} \right|\left| \overrightarrow{AC} \right|\cos \angle BAC\). 因为 \(\left| \overrightarrow{AB} \right|=\left| \overrightarrow{AC} \right|\cos \angle BAC\)，所以 \(\overrightarrow{AB}\cdot \overrightarrow{AC}=AB^2=9\). 故答案为 \(9\).
\end{solution}

\begin{problem}
已知 \(f(x)=\frac{3}{x}+2\)，则 \(f^{-1}(1)=\)\fillinblank{}.
\end{problem}

\begin{answer}
\(-3\)
\end{answer}

\begin{solution}
因为 \(f(x)=\frac{3}{x}+2\)，令 \(f(x)=1\)，即 \(\frac{3}{x}+2=1\)，解得 \(x=-3\)，故 \(f^{-1}(1)=-3\). 故答案为 \(-3\).
\end{solution}

\begin{problem}
已知二项式 \((x+a)^5\) 展开式中，\(x^2\) 的系数为 \(80\)，则 \(a=\)\fillinblank{}.
\end{problem}

\begin{answer}
\(2\)
\end{answer}

\begin{solution}
\((x+a)^5\) 的展开式的通项公式为 \(T_{r+1}=C_5^r x^{5-r}a^r\)，所以 \(x^2\) 的系数为 \(C_5^3 a^3=80\)，解得 \(a=2\). 故答案为 \(2\).
\end{solution}

\begin{problem}
已知 \(x\le 3\)，\(2x-y-2\ge 0\)，\(3x+y-8\ge 0\)，\(z=x-y\)，则 \(z\) 的最大值为\fillinblank{}.
\end{problem}

\begin{answer}
\(4\)
\end{answer}

\begin{solution}
绘制不等式组表示的平面区域如图所示，

\solutionfigure{\bitmapfigure[max width=2.8cm]{img/2021/shanghai/q7_fig1.png}}

目标函数即 \(y=x-z\)，其中 \(z\) 取得最大值时，其几何意义表示直线系在 \(y\) 轴上的截距的相反数. 据此结合目标函数的几何意义可知，目标函数在点 B 处取得最大值. 联立直线方程 \(\begin{cases}
x=3,\\
3x+y-8=0,
\end{cases}\) 可得点 B 的坐标为 \((3,-1)\). 据此可知目标函数的最大值为 \(z_{\max}=3-\left( -1 \right)=4\). 故答案为 \(4\).
\end{solution}

\begin{problem}
已知 \(\{a_n\}\) 为无穷等比数列，\(a_1=3\)，\(a_n\) 的各项和为 \(9\)，\(b_n=a_{2n}\)，则数列 \(\{b_n\}\) 的各项和为\fillinblank{}.
\end{problem}

\begin{answer}
\(\dfrac{18}{5}\)
\end{answer}

\begin{solution}
设 \(\{a_n\}\) 的公比为 \(q\)，由 \(a_1=3\)，\(a_n\) 的各项和为 \(9\)，可得 \(\frac{3}{1-q}=9\)，解得 \(q=\frac23\). 所以 \(a_n=3\times \left( \frac23 \right)^{n-1}\)，\(b_n=a_{2n}=3\times \left( \frac23 \right)^{2n-1}\). 可得数列 \(\{b_n\}\) 是首项为 \(2\)，公比为 \(\frac49\) 的等比数列，\(\sum_{n=1}^{\infty} b_n=\frac{2}{1-\frac49}=\frac{18}{5}\). 故答案为 \(\frac{18}{5}\).
\end{solution}

\begin{problem}
已知圆柱的底面圆半径为 \(1\)，高为 \(2\)，\(AB\) 为上底面圆的一条直径，C 是下底面圆周上的一个动点，则 \(\triangle ABC\) 的面积的取值范围为\fillinblank{}.
\end{problem}

\begin{answer}
\([2,\sqrt5]\)
\end{answer}

\begin{solution}
如图，上顶面圆心记为 \(O\)，下底面圆心记为 \(O'\)，

\solutionfigure{\bitmapfigure[max width=8.6cm]{img/2021/shanghai/q9_fig1.png}}

连接 \(OC\)，过点 \(C\) 作 \(CM\perp AB\)，垂足为点 \(M\)，则 \(S_{\triangle ABC}=\frac12\times AB\times CM\). 根据题意，\(AB\) 为定值 \(2\)，所以 \(S_{\triangle ABC}\) 的大小随着 \(CM\) 的长短变化而变化.

如图 \(2\) 所示，当点 \(M\) 与点 \(O\) 重合时，\(CM=OC=\sqrt{1^2+2^2}=\sqrt5\)，此时 \(S_{\triangle ABC}=\frac12\times 2\times \sqrt5=\sqrt5\).

如图 \(3\) 所示，当点 \(M\) 与点 B 重合时，\(CM\) 取最小值 \(2\)，此时 \(S_{\triangle ABC}=\frac12\times 2\times 2=2\).

综上所述，\(S_{\triangle ABC}\) 的取值范围为 \([2,\sqrt5]\). 故答案为 \([2,\sqrt5]\).
\end{solution}

\begin{problem}
已知花博会有四个不同的场馆 \(A\)，\(B\)，\(C\)，\(D\)，甲，乙两人每人选 \(2\) 个去参观，则他们的选择中，恰有一个馆相同的概率为\fillinblank{}.
\end{problem}

\begin{answer}
\(\dfrac23\)
\end{answer}

\begin{solution}
甲选 \(2\) 个去参观，有 \(C_4^2=6\) 种，乙选 \(2\) 个去参观，有 \(C_4^2=6\) 种，共有 \(6\times 6=36\) 种. 若甲乙恰有一个馆相同，则选定确定相同的馆有 \(C_4^1=4\) 种，然后从剩余 \(3\) 个馆中选 \(2\) 个进行排列，有 \(A_3^2=6\) 种，共有 \(4\times 6=24\) 种. 则对应概率 \(P=\frac{24}{36}=\frac23\). 故答案为 \(\frac23\).
\end{solution}

\begin{problem}
已知抛物线 \(y^2=2px\)（\(p>0\)），若第一象限的 \(A\)，\(B\) 在抛物线上，焦点为 \(F\)，\(|AF|=2\)，\(|BF|=4\)，\(|AB|=3\)，求直线 \(AB\) 的斜率为\fillinblank{}.
\end{problem}

\begin{answer}
\(\dfrac{\sqrt5}{2}\)
\end{answer}

\begin{solution}
如图所示，设抛物线的准线为 \(l\)，作 \(AC\perp l\) 于点 \(C\)，\(BD\perp l\) 于点 \(D\)，\(AE\perp BD\) 于点 \(E\)，

\solutionfigure{\bitmapfigure[max width=5.2cm]{img/2021/shanghai/q11_fig1.png}}

由抛物线的定义，可得 \(AC=AF=2\)，\(BD=BF=4\). 所以 \(BE=BD-AC=4-2=2\)，\(AE=\sqrt{AB^2-BE^2}=\sqrt{9-4}=\sqrt5\). 所以直线 \(AB\) 的斜率 \(k_{AB}=\tan \angle ABE=\frac{AE}{BE}=\frac{\sqrt5}{2}\). 故答案为 \(\frac{\sqrt5}{2}\).
\end{solution}

\begin{problem}
已知 \(a_i\in \N^*\)（\(i=1\)，\(2\)，\(\dots\)，\(9\)），对任意的 \(k\in \N^*\)（\(2\le k\le 8\)），\(a_k=a_{k-1}+1\) 或 \(a_k=a_{k+1}-1\) 中有且仅有一个成立，\(a_1=6\)，\(a_9=9\)，则 \(a_1+\cdots+a_9\) 的最小值为\fillinblank{}.
\end{problem}

\begin{answer}
31
\end{answer}

\begin{solution}
设 \(b_k=a_{k+1}-a_k\)，由题意可得，\(b_k\)，\(b_{k-1}\) 恰有一个为 \(1\).

如果 \(b_1=b_3=b_5=b_7=1\)，那么 \(a_1=6\)，\(a_2=7\)，\(a_3\ge 1\)，\(a_4=a_3+1\ge 2\). 同样也有，\(a_5\ge 1\)，\(a_6=a_5+1\ge 2\)，\(a_7\ge 1\)，\(a_8=a_7+1\ge 2\)，全部加起来至少是 \(6+7+1+2+1+2+1+2+9=31\).

如果 \(b_2=b_4=b_6=b_8=1\)，那么 \(a_8=8\)，\(a_2\ge 1\)，\(a_3=a_2+1\ge 2\). 同样也有，\(a_4\ge 1\)，\(a_5\ge 2\)，\(a_6\ge 1\)，\(a_7\ge 2\)，全部加起来至少是 \(6+1+2+1+2+1+2+8+9=32\).

综上所述，最小应该是 \(31\). 故答案为 \(31\).
\end{solution}
\section{单选题}

\begin{problem}
以下哪个函数既是奇函数，又是减函数（\quad）

\choices
    {\(y=-3x\)}
    {\(y=x^3\)}
    {\(y=\log_3 x\)}
    {\(y=3^x\)}
\end{problem}

\begin{answer}
A
\end{answer}

\begin{solution}
\(y=-3x\) 在 \(\R\) 上单调递减且为奇函数，\(A\) 符合题意.
因为 \(y=x^3\) 在 \(\R\) 上是增函数，\(B\) 不符合题意.
\(y=\log_3 x\)，\(y=3^x\) 为非奇非偶函数，\(C\) 不符合题意.
故选 A.
\end{solution}

\begin{problem}
已知参数方程 \(x=3t-4t^3\)，\(y=2t\sqrt{1-t^2}\)，\(t\in[-1,1]\)，以下哪个图符合该方程（\quad）

\choices
    {\choicebitmap[width=0.15\paperwidth]{img/2021/shanghai/q14_fig1.png}}
    {\choicebitmap[width=0.15\paperwidth]{img/2021/shanghai/q14_fig2.png}}
    {\choicebitmap[width=0.15\paperwidth]{img/2021/shanghai/q14_fig3.png}}
    {\choicebitmap[width=0.15\paperwidth]{img/2021/shanghai/q14_fig4.png}}
\end{problem}

\begin{answer}
B
\end{answer}

\begin{solution}
利用特殊值法进行排除. 当 \(y=0\) 时，\(t=0\)，\(1\)，\(-1\). 当 \(t=0\) 时，\(x=0\)； 当 \(t=1\) 时，\(x=-1\)； 当 \(t=-1\) 时，\(x=1\). 故当 \(y=0\) 时，\(x=0\) 或 \(1\) 或 \(-1\)，即图象经过 \((-1,0)\)，\((0,0)\)，\((1,0)\) 三个点. 对照四个选项中的图象，只有选项 B 符合要求. 故选 B.
\end{solution}

\begin{problem}
已知 \(f(x)=3\sin x+2\)，对任意的 \(x_1\in\left[ 0,\frac{\pi}{2} \right]\)，都存在 \(x_2\in\left[ 0,\frac{\pi}{2} \right]\)，使得 \(f(x_1)=2f(x_2+\theta)+2\) 成立，则下列选项中，\(\theta\) 可能的值是（\quad）

\choices
    {\(\dfrac{3\pi}{5}\)}
    {\(\dfrac{4\pi}{5}\)}
    {\(\dfrac{6\pi}{5}\)}
    {\(\dfrac{7\pi}{5}\)}
\end{problem}

\begin{answer}
B
\end{answer}

\begin{solution}
因为 \(x_1\in\left[ 0,\frac{\pi}{2} \right]\)，所以 \(\sin x_1\in[0,1]\)，从而 \(f(x_1)\in[2,5]\). 由题意，对任意的 \(x_1\in\left[ 0,\frac{\pi}{2} \right]\)，都存在 \(x_2\in\left[ 0,\frac{\pi}{2} \right]\)，使得 \(f(x_1)=2f(x_2+\theta)+2\) 成立，因此 \(f(x_2+\theta)_{\min}\le 0\)，\(f(x_2+\theta)_{\max}\ge \frac32\). 又 \(f(x)=3\sin x+2\)，所以 \(\sin(x_2+\theta)_{\min}\le -\frac23\)，\(\sin(x_2+\theta)_{\max}\ge -\frac16\). \(y=\sin x\) 在 \(x\in\left[ \frac{\pi}{2},\frac{3\pi}{2} \right]\) 上单调递减.

当 \(\theta=\frac{3\pi}{5}\) 时，\(x_2+\theta\in\left[ \frac{3\pi}{5},\frac{11\pi}{10} \right]\)，所以 \(\sin(x_2+\theta)_{\min}=\sin\frac{11\pi}{10}>\sin\frac{7\pi}{6}=-\frac12\)，
故 A 选项错误.

当 \(\theta=\frac{4\pi}{5}\) 时，\(x_2+\theta\in\left[ \frac{4\pi}{5},\frac{13\pi}{10} \right]\)，所以 \(\sin(x_2+\theta)_{\min}=\sin\frac{13\pi}{10}<\sin\frac{5\pi}{4}=-\frac{\sqrt2}{2}<-\frac23\)，\(\sin(x_2+\theta)_{\max}=\sin\frac{4\pi}{5}>0\)，
故 B 选项正确.

当 \(\theta=\frac{6\pi}{5}\) 时，\(x_2+\theta\in\left[ \frac{6\pi}{5},\frac{17\pi}{10} \right]\)，所以 \(\sin(x_2+\theta)_{\max}=\sin\frac{6\pi}{5}<\sin\frac{13\pi}{12}=-\frac{\sqrt2-\sqrt6}{4}<-\frac16\)，
故 C 选项错误.

当 \(\theta=\frac{7\pi}{5}\) 时，\(x_2+\theta\in\left[ \frac{7\pi}{5},\frac{19\pi}{10} \right]\)，所以 \(\sin(x_2+\theta)_{\max}=\sin\frac{19\pi}{10}<\sin\frac{23\pi}{12}=-\frac{\sqrt2-\sqrt6}{4}<-\frac16\)，
故 D 选项错误.

故选 B.
\end{solution}

\begin{problem}
已知两两不相等的 \(x_1\)，\(y_1\)，\(x_2\)，\(y_2\)，\(x_3\)，\(y_3\)，同时满足 \circled{1}\(x_1<y_1\)，\(x_2<y_2\)，\(x_3<y_3\)；\circled{2}\(x_1+y_1=x_2+y_2=x_3+y_3\)；\circled{3}\(x_1y_1+x_3y_3=2x_2y_2\)，以下哪个选项恒成立（\quad）

\choices
    {\(2x_2<x_1+x_3\)}
    {\(2x_2>x_1+x_3\)}
    {\(x_2^2<x_1x_3\)}
    {\(x_2^2>x_1x_3\)}
\end{problem}

\begin{answer}
A
\end{answer}

\begin{solution}
设 \(x_1+y_1=x_2+y_2=x_3+y_3=2m\)，\(x_1=m-a\)，\(y_1=m+a\)，\(x_2=m-b\)，\(y_2=m+b\)，\(x_3=m-c\)，\(y_3=m+c\). 根据题意，\(a\)，\(b\)，\(c>0\)，且 \(a^2+c^2=2b^2\). 因而
\(x_1+x_3-2x_2=\left( m-a \right)+\left( m-c \right)-2\left( m-b \right)=2b-\left( a+c \right)\). 因为 \(\left( 2b \right)^2-\left( a+c \right)^2=2\left( a^2+c^2 \right)-\left( a+c \right)^2=(a-c)^2>0\)，所以 \(x_1+x_3-2x_2=2b-\left( a+c \right)>0\). 所以 A 项正确，B 错误. \(x_1x_3-x_2^2=\left( m-a \right)\left( m-c \right)-\left( m-b \right)^2=\left( 2b-a-c \right)m-\frac{\left( a-c \right)^2}{2}\). 而上面已证 \(\left( 2b-a-c \right)>0\)，因为不知道 \(m\) 的正负，所以该式子的正负无法恒定. 故选 A.
\end{solution}
\section{解答题}

\begin{problem}
如图，在长方体 \(ABCD-A_1B_1C_1D_1\) 中，已知 \(AB=BC=2\)，\(AA_1=3\).

\begin{enumerate}
\item 若 \(P\) 是棱 \(A_1D_1\) 上的动点，求三棱锥 \(C-PAD\) 的体积；
\item 求直线 \(AB_1\) 与平面 \(ACC_1A_1\) 的夹角大小.
\end{enumerate}

\bitmapfigure[max width=5.2cm]{img/2021/shanghai/q17_fig1.png}
\end{problem}

\begin{solution}
\begin{enumerate}
\item 因为平面 \(PAD\) 就是长方体的左侧面，点 \(C\) 到平面 \(PAD\) 的距离为 \(BC=2\). 又 \(S_{\triangle PAD}=\frac12\times AD\times AA_1=\frac12\times 2\times 3=3\)，所以 \(V_{C-PAD}=\frac13 S_{\triangle PAD}\times 2=\frac13\times 3\times 2=2\).

\item

连接 \(A_1C_1\) 与 \(B_1D_1\)，交于点 \(O\). 因为 \(AB=BC\)，所以四边形 \(A_1B_1C_1D_1\) 为正方形，则 \(OB_1\perp OA_1\). 又 \(AA_1\perp OB_1\)，\(OA_1\cap AA_1=A_1\)，所以 \(OB_1\perp\) 平面 \(ACC_1A_1\). 于是直线 \(AB_1\) 与平面 \(ACC_1A_1\) 所成的角为 \(\angle OAB_1\)，且 \(\sin \angle OAB_1=\frac{OB_1}{AB_1}=\frac{\frac{\sqrt{2^2+2^2}}{2}}{\sqrt{2^2+3^2}}=\frac{\sqrt{26}}{13}\). 所以直线 \(AB_1\) 与平面 \(ACC_1A_1\) 所成的角为 \(\arcsin \frac{\sqrt{26}}{13}\).
\end{enumerate}
\solutionfigure{\bitmapfigure[max width=5.2cm]{img/2021/shanghai/q17_solution_fig1.png}}
\end{solution}

\begin{problem}
在 \(\triangle ABC\) 中，已知 \(a=3\)，\(b=2c\).

\begin{enumerate}
\item 若 \(A=\frac{2\pi}{3}\)，求 \(S_{\triangle ABC}\)；
\item 若 \(2\sin B-\sin C=1\)，求 \(C_{\triangle ABC}\).
\end{enumerate}
\end{problem}

\begin{solution}
\begin{enumerate}
\item 由余弦定理得 \(\cos A=-\frac12=\frac{b^2+c^2-a^2}{2bc}=\frac{5c^2-9}{4c^2}\)，解得 \(c^2=\frac97\). 所以 \(S_{\triangle ABC}=\frac12bc\sin A=\frac{\sqrt3}{4}\times 2c^2=\frac{9\sqrt3}{14}\).

\item 因为 \(b=2c\)，所以由正弦定理得 \(\sin B=2\sin C\). 又 \(2\sin B-\sin C=1\)，所以 \(\sin C=\frac13\)，\(\sin B=\frac23\). 由 \(b=2c\) 及正弦定理可知 \(\sin B=2\sin C\)，且因边 \(b>c\)，有 \(B>C\). 又 \(B+C<\pi\)，故 \(C<\frac\pi2\)，从而 \(\cos C=\sqrt{1-\left( \frac13 \right)^2}=\frac{2\sqrt2}{3}\). 由余弦定理得 \(c^2=a^2+b^2-2ab\cos C\). 又 \(a=3\)，\(b=2c\)，所以 \(c^2=9+4c^2-8\sqrt2\,c\)，即 \(3c^2-8\sqrt2\,c+9=0\)，解得 \(c=\frac{4\sqrt2\pm \sqrt5}{3}\).

当 \(c=\frac{4\sqrt2+\sqrt5}{3}\) 时，\(b=\frac{8\sqrt2+2\sqrt5}{3}\)，所以 \(C_{\triangle ABC}=3+4\sqrt2+\sqrt5\).

当 \(c=\frac{4\sqrt2-\sqrt5}{3}\) 时，\(b=\frac{8\sqrt2-2\sqrt5}{3}\)，所以 \(C_{\triangle ABC}=3+4\sqrt2-\sqrt5\).
\end{enumerate}
\end{solution}

\begin{problem}
已知一企业今年第一季度的营业额为 \(1.1\) 亿元，往后每个季度增加 \(0.05\) 亿元，第一季度的利润为 \(0.16\) 亿元，往后每一季度比前一季度增长 \(4\%\).

\begin{enumerate}
\item 求今年起的前 \(20\) 个季度的总营业额；
\item 请问哪一季度的利润首次超过该季度营业额的 \(18\%\)？
\end{enumerate}
\end{problem}

\begin{solution}
\begin{enumerate}
\item 由题意可知，可将每个季度的营业额看作等差数列，则首项 \(a_1=1.1\)，公差 \(d=0.05\). 所以 \(S_{20}=20a_1+\frac{20\left( 20-1 \right)}{2}d=20\times 1.1+10\times 19\times 0.05=31.5\). 即营业额前 \(20\) 个季度的和为 \(31.5\) 亿元.

\item 【法一】假设今年第一季度往后的第 \(n\)（\(n\in \N^*\)）季度的利润首次超过该季度营业额的 \(18\%\)，则 \(0.16\times \left( 1+4\% \right)^n>\left( 1.1+0.05n \right)\cdot 18\%\). 令 \(f(n)=0.16\times \left( 1+4\% \right)^n-\left( 1.1+0.05n \right)\cdot 18\%\)，\(n\in \N^*\)，即要解 \(f(n)>0\). 则当 \(n\ge 2\) 时，\(f(n)-f(n-1)=0.0064\cdot \left( 1+4\% \right)^{n-1}-0.009\). 令 \(f(n)-f(n-1)>0\)，解得 \(n\ge 10\)，即当 \(1\le n\le 9\) 时，\(f(n)\) 递减；当 \(n\ge 10\) 时，\(f(n)\) 递增. 由于 \(f(1)<0\)，因此 \(f(n)>0\) 的解只能在 \(n\ge 10\) 时取得. 经检验，\(f(24)<0\)，\(f(25)>0\)，所以今年第一季度往后的第 \(25\) 个季度的利润首次超过该季度营业额的 \(18\%\).

【法二】设今年第一季度往后的第 \(n\)（\(n\in \N^*\)）季度的利润与该季度营业额的比为 \(a_n\)，则 \(\frac{a_{n+1}}{a_n}=\frac{1.04\left( 1.05+0.05n \right)}{1.1+0.05n}=1.04-\frac{1.04}{22+n}=1+0.04\left( 1-\frac{26}{22+n} \right)\)，所以数列 \(\{a_n\}\) 满足 \(a_1>a_2>a_3>a_4=a_5<a_6<a_7<\cdots\). 注意到，\(a_{25}=0.178\cdots\)，\(a_{26}=0.181\cdots\)，所以今年第一季度往后的第 \(25\) 个季度利润首次超过该季度营业额的 \(18\%\).
\end{enumerate}
\end{solution}

\begin{problem}
已知椭圆 \(\Gamma:\frac{x^2}{2}+y^2=1\)，\(F_1\)，\(F_2\) 是其左，右焦点，直线 \(l\) 过点 \(P(m,0)\)（\(m\le -\sqrt2\)），交椭圆于 \(A\)，\(B\) 两点，且 \(A\)，\(B\) 在 \(x\) 轴上方，点 \(A\) 在线段 \(BP\) 上.

\begin{enumerate}
\item 若 \(B\) 是上顶点，\(\left| \overrightarrow{BF_1} \right|=\left| \overrightarrow{PF_1} \right|\)，求 \(m\) 的值；
\item 若 \(\overrightarrow{F_1A}\cdot \overrightarrow{F_2A}=\frac13\)，且原点 \(O\) 到直线 \(l\) 的距离为 \(\frac{4\sqrt{15}}{15}\)，求直线 \(l\) 的方程；
\item 证明：对于任意 \(m<-\sqrt2\)，使得 \(\overrightarrow{F_1A}\parallel \overrightarrow{F_2B}\) 的直线有且仅有一条.
\end{enumerate}
\end{problem}

\begin{solution}
\begin{enumerate}
\item 因为 \(\Gamma\) 的方程为 \(\frac{x^2}{2}+y^2=1\)，所以 \(a^2=2\)，\(b^2=1\)，从而 \(c^2=a^2-b^2=1\). 所以 \(F_1(-1,0)\)，\(F_2(1,0)\). 若 \(B\) 为 \(\Gamma\) 的上顶点，则 \(B(0,1)\)，所以 \(\left| BF_1 \right|=\sqrt2\)，\(\left| PF_1 \right|=-1-m\). 又 \(\left| BF_1 \right|=\left| PF_1 \right|\)，所以 \(m=-1-\sqrt2\).

\item 设点 \(A\left( \sqrt2\cos\theta,\sin\theta \right)\)，则 \(\overrightarrow{F_1A}\cdot \overrightarrow{F_2A}=\left( \sqrt2\cos\theta+1 \right)\left( \sqrt2\cos\theta-1 \right)+\sin^2\theta=2\cos^2\theta-1+\sin^2\theta=\frac13\). 因为 \(A\) 在线段 \(BP\) 上，横坐标小于 \(0\)，解得 \(\cos\theta=-\frac{\sqrt3}{3}\). 故 \(A\left( -\frac{\sqrt6}{3},\frac{\sqrt6}{3} \right)\). 设直线 \(l\) 的方程为 \(y=kx+\frac{\sqrt6}{3}k+\frac{\sqrt6}{3}\)（\(k>0\)）. 由原点 \(O\) 到直线 \(l\) 的距离为 \(\frac{4\sqrt{15}}{15}\)，\(\frac{\left| \frac{\sqrt6}{3}k+\frac{\sqrt6}{3} \right|}{\sqrt{1+k^2}}=\frac{4\sqrt{15}}{15}\)，化简可得 \(3k^2-10k+3=0\)，解得 \(k=3\) 或 \(k=\frac13\). 故直线 \(l\) 的方程为 \(y=3x+\frac{4\sqrt6}{3}\) 或 \(y=\frac13x+\frac{4\sqrt6}{9}\). 其中 \(y=3x+\frac{4\sqrt6}{3}\) 舍去，无法满足 \(m<-\sqrt2\). 所以直线 \(l\) 的方程为 \(y=\frac13x+\frac{4\sqrt6}{9}\).

\item 联立方程组
\examdisplaycases{}{
y=kx-km,\\
\frac{x^2}{2}+y^2=1,
}
可得 \(\left( 1+2k^2 \right)x^2-4k^2mx+2k^2m^2-2=0\). 设 \(A(x_1,y_1)\)，\(B(x_2,y_2)\)，则 \(x_1+x_2=\frac{4k^2m}{1+2k^2}\)，\(x_1x_2=\frac{2k^2m^2-2}{1+2k^2}\). 因为 \(\overrightarrow{F_1A}\parallel \overrightarrow{F_2B}\)，所以 \(\left( x_2-1 \right)y_1=\left( x_1+1 \right)y_2\). 又 \(y=kx-km\)，故化简为 \(x_1-x_2=-\frac{2}{1+2k^2}\). 又
\[
\left| x_1-x_2 \right|
=\sqrt{\left( x_1+x_2 \right)^2-4x_1x_2}
=\frac{\sqrt{16k^2-8k^2m^2+8}}{1+2k^2}
=\left| -\frac{2}{1+2k^2} \right|
\]
两边同时平方可得 \(4k^2-2k^2m^2+1=0\)，整理可得 \(k^2=-\frac{1}{4-2m^2}\). 当 \(m<-\sqrt2\) 时，\(k^2=-\frac{1}{4-2m^2}>0\). 因为点 \(A\)，\(B\) 在 \(x\) 轴上方，所以 \(k\) 有且仅有一个解. 故对于任意 \(m<-\sqrt2\)，使得 \(\overrightarrow{F_1A}\parallel \overrightarrow{F_2B}\) 的直线有且仅有一条.
\end{enumerate}
\end{solution}

\begin{problem}
已知 \(x_1\)，\(x_2\in \R\)，若对任意的 \(x_2-x_1\in S\)，\(f(x_2)-f(x_1)\in S\)，则有定义：\(f(x)\) 是在 \(S\) 关联的.

\begin{enumerate}
\item 判断和证明 \(f(x)=2x-1\) 是否在 \([0,+\infty)\) 关联？是否有 \([0,1]\) 关联？
\item 若 \(f(x)\) 是在 \(\{3\}\) 关联的，\(f(x)\) 在 \(x\in[0,3)\) 时，\(f(x)=x^2-2x\)，求解不等式 \(2\le f(x)\le 3\)；
\item 证明：\(f(x)\) 是 \(\{1\}\) 关联的，且是在 \([0,+\infty)\) 关联的，当且仅当“\(f(x)\) 在 \([1,2]\) 是关联的”.
\end{enumerate}
\end{problem}

\begin{solution}
\begin{enumerate}
\item \(f(x)\) 在 \([0,+\infty)\) 关联，在 \([0,1]\) 不关联. 任取 \(x_1-x_2\in[0,+\infty)\)，则 \(f(x_1)-f(x_2)=2\left( x_1-x_2 \right)\in[0,+\infty)\). 所以 \(f(x)\) 在 \([0,+\infty)\) 关联.

取 \(x_1=1\)，\(x_2=0\)，则 \(x_1-x_2=1\in[0,1]\)，\(f(x_1)-f(x_2)=2\left( x_1-x_2 \right)=2\notin[0,1]\). 所以 \(f(x)\) 在 \([0,1]\) 不关联.

\item 因为 \(f(x)\) 在 \(\{3\}\) 关联，所以对于任意 \(x_1-x_2=3\)，都有 \(f(x_1)-f(x_2)=3\). 所以对任意 \(x\)，都有 \(f(x+3)-f(x)=3\). 由 \(x\in[0,3)\) 时，\(f(x)=x^2-2x\)，得 \(f(x)\) 在 \(x\in[0,3)\) 的值域为 \([-1,3)\). 所以 \(f(x)\) 在 \(x\in[3,6)\) 的值域为 \([2,6)\). 因此 \(2\le f(x)\le 3\) 仅在 \(x\in[0,3)\) 或 \(x\in[3,6)\) 上有解.

当 \(x\in[0,3)\) 时，\(f(x)=x^2-2x\). 令 \(2\le x^2-2x\le 3\)，解得 \(\sqrt3+1\le x<3\).

当 \(x\in[3,6)\) 时，\(f(x)=f(x-3)+3=x^2-8x+18\). 令 \(2\le x^2-8x+18\le 3\)，解得 \(3\le x\le 5\).

所以不等式 \(2\le f(x)\le 3\) 的解为 \([\sqrt3+1,5]\).

\item 先证明：\(f(x)\) 是在 \(\{1\}\) 关联的，且是在 \([0,+\infty)\) 关联的 \(\Rightarrow f(x)\) 在 \([1,2]\) 是关联的. 由已知条件可得 \(f(x+1)=f(x)+1\)，所以 \(f(x+n)=f(x)+n\)，\(n\in \Z\). 又因为 \(f(x)\) 是在 \([0,+\infty)\) 关联的，所以任意 \(x_2>x_1\)，都有 \(f(x_2)>f(x_1)\). 若 \(1\le x_2-x_1\le 2\)，则 \(x_1+1\le x_2\le x_1+2\)，从而 \(f(x_1+1)\le f(x_2)\le f(x_1+2)\)，即 \(f(x_1)+1\le f(x_2)\le f(x_1)+2\)，所以 \(1\le f(x_2)-f(x_1)\le 2\). 因此 \(f(x)\) 是 \([1,2]\) 关联.

再证明：\(f(x)\) 在 \([1,2]\) 是关联的 \(\Rightarrow f(x)\) 是在 \(\{1\}\) 关联的，且是在 \([0,+\infty)\) 关联的. 因为 \(f(x)\) 在 \([1,2]\) 是关联的，所以任取 \(x_1-x_2\in[1,2]\)，都有 \(f(x_1)-f(x_2)\in[1,2]\). 下面用反证法证明 \(f(x+1)-f(x)=1\). 若 \(f(x+1)-f(x)>1\)，则 \(f(x+2)-f(x)=f(x+2)-f(x+1)+f(x+1)-f(x)>2\)，这与 \(f(x)\) 在 \([1,2]\) 是关联的矛盾. 若 \(f(x+1)-f(x)<1\)，而 \(f(x)\) 在 \([1,2]\) 是关联的，则 \(f(x+1)-f(x)\ge 1\)，矛盾. 所以 \(f(x+1)-f(x)=1\) 成立，即 \(f(x)\) 是在 \(\{1\}\) 关联的. 再证明 \(f(x)\) 是在 \([0,+\infty)\) 关联的. 任取 \(x_1-x_2\in[n,+\infty)\)（\(n\in \N\)），则存在 \(n\in \N\)，使得 \(x_1-x_2\in[n,n+1]\). 于是 \(1\le x_1-(n-1)-x_2\le 2\). 所以 \(f\left[ x_1-(n-1) \right]-f(x_2)=f(x_1)-(n-1)-f(x_2)\in[1,2]\). 从而 \(f(x_1)-f(x_2)\in[n,n+1]\subset[0,+\infty)\). 因此 \(f(x)\) 是在 \([0,+\infty)\) 关联的.

综上所述，\(f(x)\) 是 \(\{1\}\) 关联的，且是在 \([0,+\infty)\) 关联的，当且仅当“\(f(x)\) 在 \([1,2]\) 是关联的”.
\end{enumerate}
\end{solution}
